\newpage
\section{Planteamiento del Problema}

En el contexto de la Facultad de Ciencias y Tecnología de la Universidad Mayor de San Simón, específicamente en las carreras de Ingeniería en Informática e Ingeniería de Sistemas, la formación en el área de seguridad informática forma parte del proceso académico en niveles avanzados mediante la asignatura de Seguridad de Sistemas. Esta materia constituye una de las principales instancias de abordaje formal de la ciberseguridad dentro del plan de estudios, por lo que resulta relevante para el desarrollo de competencias técnicas relacionadas con la identificación, análisis y respuesta ante amenazas digitales.

Sin embargo, la situación actual evidencia que el desarrollo de esta asignatura mantiene un enfoque predominantemente teórico. Los contenidos son expuestos principalmente por el docente, mientras que las actividades prácticas suelen ser realizadas de forma autónoma por los estudiantes fuera del aula, generalmente mediante el uso de herramientas como máquinas virtuales, laboratorios individuales o ejercicios complementarios. Esta dinámica reduce las posibilidades de experimentación guiada, acompañamiento continuo y retroalimentación inmediata durante el proceso de aprendizaje.

Asimismo, se observa una elevada cantidad de estudiantes inscritos en la asignatura, alcanzando aproximadamente 130 estudiantes en la gestión 2026-I y alrededor de 180 estudiantes en la gestión 2025-II. Esta cantidad de participantes dificulta el seguimiento individualizado, la supervisión de prácticas y la corrección oportuna de errores durante el desarrollo de ejercicios técnicos. De igual manera, los mecanismos de evaluación se centran principalmente en exámenes teóricos y trabajos expositivos, lo cual refuerza el carácter conceptual de la materia y limita la valoración de habilidades prácticas aplicadas.

La ciberseguridad, por su propia naturaleza, requiere que el estudiante no solo comprenda conceptos, sino que también practique la toma de decisiones frente a escenarios de riesgo, identifique vulnerabilidades, analice comportamientos sospechosos y aplique respuestas adecuadas ante incidentes. Por ello, cuando el aprendizaje se desarrolla con poca práctica guiada, se genera una brecha entre el conocimiento teórico adquirido y la capacidad real de aplicarlo en situaciones similares a las que pueden presentarse en un entorno profesional.

Frente a esta situación, los simuladores y las estrategias de gamificación representan una alternativa pedagógica pertinente, ya que permiten recrear escenarios controlados de amenazas digitales, incorporar retos progresivos, ofrecer retroalimentación inmediata y mantener la participación activa del estudiante. A diferencia de una práctica aislada, un entorno gamificado puede organizar el aprendizaje mediante misiones, niveles, objetivos y recompensas, facilitando que el estudiante experimente, se equivoque, corrija y fortalezca sus habilidades de manera progresiva.

En este sentido, se evidencia la necesidad de desarrollar una herramienta educativa basada en simulación gamificada que apoye el aprendizaje práctico de la ciberseguridad en estudiantes de Ingeniería en Informática e Ingeniería de Sistemas de la Universidad Mayor de San Simón. Esta herramienta permitiría complementar el enfoque teórico de la asignatura mediante escenarios interactivos orientados al entrenamiento de habilidades prácticas, contribuyendo a reducir la brecha existente entre la enseñanza conceptual y la aplicación técnica de conocimientos en ciberseguridad.

\subsection*{Formulación del problema}
\addcontentsline{toc}{subsection}{Formulación del problema}

\noindent ¿De qué manera el desarrollo de un videojuego de simulación gamificada —basado en
escenarios de amenazas digitales y retroalimentación inmediata— contribuye al fortalecimiento
de las habilidades prácticas de respuesta ante incidentes en estudiantes de Ingeniería
informática y sistemas de la UMSS?